
\section{用\trDistinguishedAffineBase{}刻画\trQC{}性}

因为\trQCSheaf{}在局部上都呈一种十分特殊的形态，要想「认识」一个\trQCSheaf，我们只需知道诸\trSection{}在每个\trAffine\trOpenSet{}上是什么，以及如何把\trSection{}从一个\trAffine\trOpenSet{} $\mathscr{U}$ 限制到 $\mathscr{U}$ 的特异\trAffine\trOpenSubset{}上。我们将通过定义扎里斯基\trTopology{}的所谓\textbf{\trDistinguishedAffineBase{}}（distinguished affine base）来把这件事讲精确——它并非通常意义下的\trBase。这是对「\trSheafOnBase」观念的一种精细化。这一番讨论的要点在于给出\trQC{}性的一个有用刻画，不过你或许愿意径直跳到 \S6.2.3。

\trDistinguishedAffineBase{}中的\trOpenSet，就是 $X$ 的诸\trAffine\trOpenSubset。我们已经观察到它们构成一个\trBase。但请把这一事实忘掉。我们喜爱的是诸\trDistinguishedOpen{} $\Spec A_f\hookrightarrow\Spec A$；至于随便哪个\trAffine\trOpenSubset{}开嵌入另一个之中的那种\trOpenEmbedding，我们并无真正的理解。所以我们干脆只记住那些「好」的\trInclusionMap。

\textbf{6.2.1. \trhDefinition{}。}\quad 概形 $X$ 的\textbf{\trDistinguishedAffineBase{}}，就是「诸\trAffine\trOpenSet{}与诸特异\trInclusionMap」这份数据。

换言之，我们只记住\trOpenSet{}中的一部分（诸\trAffine\trOpenSet），而且\emph{只记住它们之间的一部分\trMorphism}（诸特异\trMorphism）。给行家的：若你把一个\trTopology{}想成一个\trCategory{}（即\trOpenSet{}的\trCategory），那么我们所描述的便是一个\trSubcategory。

我们可以按显而易见的方式定义\trDistinguishedAffineBase{}上的\trSheaf{}：给每个\trAffine\trOpenSet{}配一个\trSet{}（或\trAbelianGroups、或\trRing），并且我们知道如何限制到诸\trDistinguishedOpen{}上。（\trIdentityAxiom{}与\trGluability{}公理的陈述，你应当自行想通。）

给定 $X$ 上的一个\trSheaf{} $\mathscr{F}$，我们便得到\trDistinguishedAffineBase{}上的一个\trSheaf。你能猜到我们正往哪里去：我们将证明，\trSheaf{}的全部信息都蕴含于\trDistinguishedAffineBase{}上的\trSheaf{}的信息之中。

作为热身，我们可以如下地复原诸\trStalk{}。（我们将隐含地只用到下面这个事实。设有\trOpenSubset{}组成的一个族，以及这些子集之间的\emph{某些}\trInclusionMap，使得对任取的一点 $p\in\mathscr{U}\cap\mathscr{V}$（$\mathscr{U}$、$\mathscr{V}$ 都属于我们的\trOpenSubset{}族），总存在某个 $\mathscr{W}$ 包含 $p$ 且含于 $\mathscr{U}$ 与 $\mathscr{V}$ 之内，而 $\mathscr{W}\hookrightarrow\mathscr{U}$ 与 $\mathscr{W}\hookrightarrow\mathscr{V}$ 都属于我们的\trInclusionMap{}之族。在我们此处考虑的情形，这正是关键的命题~5.3.1：给定 $X$ 中任意两个\trAffine\trOpenSet{} $\Spec A$、$\Spec B$，$\Spec A\cap\Spec B$ 能够被一族「在 $\Spec A$ 与 $\Spec B$ 中同时特异」的\trAffine\trOpenSet{}覆盖。用华丽的语言说：由\trAffine\trOpenSet{}与特异\trInclusionMap{}组成的\trCategory{}构成一个\trFiltered{}集。）

\trStalk{} $\mathscr{F}_p$ 是\trColimit{} $\colim(f\in\mathscr{F}(\mathscr{U}))$，该\trColimit{}取遍所有含 $p$ 的\trOpenSet。我们把它与 $\colim(f\in\mathscr{F}(\mathscr{U}))$ 相比较，这里该\trColimit{}取遍所有含 $p$ 的\trAffine\trOpenSet{}与所有特异\trInclusionMap。你可以检验：其中一方的元素与另一方的元素相互对应。（仔细想想这件事！）

\textbf{6.2.A.（练习）}\quad 证明：\trDistinguishedAffineBase{}上一个\trSheaf{}的\trSection{}由该\trSection{}的诸\trGerm{}确定。

\textbf{6.2.2.\ \trhTheorem{}。}\ ——\
\begin{enumerate}[label=(\alph*)]
  \item \emph{\trDistinguishedAffineBase{}上的\trSheaf{} $\mathscr{F}^{\flat}$ 确定一个\trSheaf{} $\mathscr{F}$：它在相差唯一\trIsomorphism{}的意义下唯一，且限制到仿射基底上时即为 $\mathscr{F}^{\flat}$。（因此，若你从一个\trSheaf{}出发，取其在\trDistinguishedAffineBase{}上的\trSheaf，再取由此诱导的\trSheaf，你得到的正是出发时的那个\trSheaf。）}
  \item \emph{\trDistinguishedAffineBase{}上的\trSheaf{}之间的\trMorphism，唯一地确定诸\trSheaf{}之间的\trMorphism。}
  \item \emph{「\trDistinguishedAffineBase{}上」的 $\OO_X$-\trModule{}给出一个 $\OO_X$-\trModule。}
\end{enumerate}

这个证明与我们在 \S2.5 中那个「诸\trSheaf{}（在本质上）就是\trBase{}上的诸\trSheaf」的论证一模一样，用的也是「\trCompatibleGerm\trSheaf」这一构造。之所以不厌其烦地再讲一遍，主要是想让你亲眼看到：这里所需要的一切，无非是诸\trOpenSet{}构成一个\trFiltered{}集（或者在当前情形下，由\trOpenSet{}与特异\trInclusionMap{}构成的\trCategory{}是\trFiltered{}的）。

给行家的：(a) 与 (b) 所描述的，是「$X$ 的\trZariskiTopology{}上的诸\trSheaf」与「$X$ 的\trDistinguishedAffineBase{}上的诸\trSheaf」之间的一个\trEquivalenceOfCategories。

\emph{证明。}\ (a)~设 $\mathscr{F}^{\flat}$ 是\trDistinguishedAffineBase{}上的一个\trSheaf。那么我们可以定义诸\trStalk。

对 $X$ 的\emph{任何}\trOpenSet{} $\mathscr{U}$，定义\trCompatibleGerm\trSheaf
\[
\mathscr{F}(\mathscr{U}):=\Bigl\{(f_p\in\mathscr{F}^{\flat}_p)_{p\in\mathscr{U}}:\ \begin{array}{l}
\text{对一切 }p\in\mathscr{U}\text{，存在 }\mathscr{U}_p\text{ 使 }p\in\mathscr{U}_p\subset\mathscr{U}\text{，}\\
\text{且存在 }s\in\mathscr{F}^{\flat}(\mathscr{U}_p)\text{，使得对一切 }q\in\mathscr{U}_p\text{ 都有 }s_q=f_q
\end{array}\Bigr\},
\]
其中每个 $\mathscr{U}_p$ 都在我们的\trBase{}中，而 $s_q$ 指「$s$ 在 $q$ 处的\trGerm{}」。（照例，想为空集操心的人自便。）

这确实是一个\trSheaf{}：请说服自己，\trRestrictionMap、\trIdentityAxiom{}与\trGluability{}我们都已具备，而且来得相当轻易。

我接下来断言：若 $\mathscr{U}$ 在我们的\trBase{}中，则 $\mathscr{F}(\mathscr{U})=\mathscr{F}^{\flat}(\mathscr{U})$。我们显然有一个映射 $\mathscr{F}^{\flat}(\mathscr{U})\to\mathscr{F}(\mathscr{U})$。它在诸\trStalk{}上是\trIsomorphism，故由\cref{ex:isomorphisms-determined-by-stalks}，它是\trIsomorphism。

\textbf{6.2.B.（练习）}\quad 证明 (b)（对照\cref{ex:sheaf-morphisms-on-base}）。

\textbf{6.2.C.（练习）}\quad 证明 (c)（对照\cref{thm:ox-modules-on-base}）。

至此，定理~6.2.2 的证明便告完成。\hfill$\square$

\textbf{6.2.3. 用特异\trInclusionMap{}刻画\trQCSheaf。}\quad 我们将用这一视角，在诸 $\OO_X$-\trModule{}之中给出\trQCSheaf{}的一个有用刻画。设 $\mathscr{F}$ 是一个 $\OO_X$-\trModule，并设 $\Spec A_f\hookrightarrow\Spec A\subset X$ 是 $X$ 的一个\trAffine\trOpenSubscheme{}的\trDistinguishedOpen\trSubscheme。设 $\phi\colon\Gamma(\Spec A,\mathscr{F})\to\Gamma(\Spec A_f,\mathscr{F})$ 是\trRestrictionMap。$\phi$ 的源是一个 $A$-\trModule，目标是 $A_f$-\trModule，故由\trLocalization{}的\trUniversalProperty{}（\cref{ex:rings-localization-universal-property}），$\phi$ 自然地分解为：
\[
\begin{tikzcd}[row sep=2.6em, column sep=5.5em]
\Gamma(\Spec A,\mathscr{F}) \ar[rr,"\phi"] \ar[dr,"\otimes_A A_f"'] && \Gamma(\Spec A_f,\mathscr{F}) \\
& \Gamma(\Spec A,\mathscr{F})_f \ar[ur,"\alpha"] &
\end{tikzcd}
\]

\textbf{6.2.D.（十分重要的练习）}\quad 证明：一个 $\OO_X$-\trModule{} $\mathscr{F}$ 是\trQC{}的，当且仅当对每个这样的特异包含 $\Spec A_f\hookrightarrow\Spec A$，$\alpha$ 都是\trIsomorphism。

于是，一个\trQCSheaf{}（等价地）就是这样的数据：对每个\trAffine\trOpenSubset{}配有一个\trModule{}（即相应\trRing{}上的一个\trModule），使得\trDistinguishedOpen{} $\Spec A_f$ 上的那个\trModule，由对 $\Spec A$ 上的\trModule{}作\trLocalization{}而给出。

\textbf{6.2.E.（好练手：幂零层）}\quad 设 $A$ 是一个\trRing，$f\in A$。证明 $\mathfrak{N}(A_f)\cong\mathfrak{N}(A)_f$。用这件事在任意概形 $X$ 上定义／构造\trQC{}的\textbf{\trNilpotent\trSheaf}（sheaf of nilpotents）。这是（$\OO_X$ 的）\trIdeal\trSheaf{}的一个例子。

\textbf{6.2.4.}\ \emph{\trTensorProduct。}\quad $X$ 上\trQCSheaf{}范畴中的\trTensorProduct，也能借助\trAffine\trOpenSet{}得到一个干净的描述。

\textbf{6.2.F.（练习）}\quad 设 $\mathscr{F}$ 与 $\mathscr{G}$ 是概形 $X$ 上的两个\trQCSheaf。证明 $\mathscr{F}\otimes\mathscr{G}$ 是一个\trQCSheaf，并可由下述信息描述：若 $\Spec A\subset X$ 是一个\trAffine\trOpenSubset，且 $\Gamma(\Spec A,\mathscr{F})=M$、$\Gamma(\Spec A,\mathscr{G})=N$，则 $\Gamma(\Spec A,\mathscr{F}\otimes\mathscr{G})=M\otimes_A N$，而且\trRestrictionMap{} $\Gamma(\Spec A,\mathscr{F}\otimes\mathscr{G})\to\Gamma(\Spec A_f,\mathscr{F}\otimes\mathscr{G})$ 恰是\trLocalization{}映射 $M\otimes_A N\to(M\otimes_A N)_f\cong M_f\otimes_{A_f}N_f$。（我们用到了代数事实 $(M\otimes_A N)_f\cong M_f\otimes_{A_f}N_f$。愿意的话，你既可以用\trUniversalProperty{}证明它，也可以用显式构造证明。）

请注意，多亏了\trDistinguishedAffineBase{}背后的这套机器，\trSheafification{}已被自动照料妥当。这是一项我们将反复使用的特性：对一般的\trSheaf{}需要\trSheafification{}的那些涉及\trQCSheaf{}的构造，一旦放到\trDistinguishedAffineBase{}上考虑，就不再需要\trSheafification。单射性、满射性、\trKernel、\trTensorProduct{}等等都可以在\trAffine\trOpenSet{}上计算——连同这一点，这正是「用\trAffine\trOpenSet{}的语言思考\trQCSheaf」格外方便的原因。（在层 $\mathscr{H}om$ 的情形有一个小小的警示；见练习~14.3.A。）

\textbf{6.2.5.}\ \extrahard\ \emph{雅致的旁注，但眼下或许过于炫技。}\quad 在练习~6.2.F 中，那个\trTensorProduct{}取在\trQCSheaf{}的范畴中。但事实上，它甚至就是 $\OO_X$-\trModule{}范畴中的\trTensorProduct。为了证明「$\mathscr{F}\otimes_{\catQCo_X}\mathscr{G}$ 满足 $\mathscr{F}\otimes_{\OO_X}\mathscr{G}$ 的\trUniversalProperty{}」，请先证明下面这件事。设 $\mathscr{M}$ 是 $X$ 上的一个\trQCSheaf，$\mathscr{N}$ 是任意一个 $\OO_X$-\trModule。设 $\Spec A\subset X$ 是一个\trAffine\trOpenSubset，且 $M=\mathscr{M}(\Spec A)$。则「\trGlobalSections」映射
\[
\Hom_{\Spec A}\bigl(\mathscr{M}|_{\Spec A},\ \mathscr{N}|_{\Spec A}\bigr)\longrightarrow\Hom_A\bigl(M,\ \mathscr{N}(\Spec A)\bigr)
\]
是\trIsomorphism。我们的口号：在一个\trAffineScheme{}上，从\trQCSheaf{}到\emph{任何} $\OO$-\trModule{}的映射都由\trGlobalSections{}确定。你将在练习~7.3.G 的一句夹注评论中听到这一口号的回声。

\textbf{6.2.6. \trhDefinition{}。}\quad 给定 $\mathscr{F}$ 的一个\trSection{} $s$ 与 $\mathscr{G}$ 的一个\trSection{} $t$，我们就有 $\mathscr{F}\otimes\mathscr{G}$ 的一个\trSection{} $s\otimes t$。

\textbf{6.2.7.}\ \emph{$X_f$ 与 Qcqs 引理。}\footnote{译注：Qcqs 是 quasicompact quasiseparated（拟紧拟分离）的缩写。}\quad 下面这个结果日后将派上用场：借它可证某些构造确实给出\trQCSheaf。

\textbf{6.2.8.}\ \emph{\trhDefinition{}：$X_f$。}\quad 设 $X$ 是一个概形，$f\in\Gamma(X,\OO_X)$ 是一个函数。令 $X_f\subset X$ 为 $f$ 在其上不消失的子集。由\cref{ex:vanishing-locus-closed}(a)，$X_f$ 是开的。（我们避免使用记号 $D(f)$，以免暗示 $X$ 是仿射的。）

\textbf{6.2.9. Qcqs 引理。}\ ——\ \emph{设 $X$ 是一个\trQuasicompact{}且\trQuasiSeparated{}的概形，$\mathscr{F}$ 是 $X$ 上的一个\trQCSheaf，而 $f\in\Gamma(X,\OO_X)$ 是 $X$ 上的一个函数。则\trRestrictionMap
\[
\mathrm{res}_{X_f\subset X}\colon\ \Gamma(X,\mathscr{F})\longrightarrow\Gamma(X_f,\mathscr{F})
\]
恰是\trLocalization{}：存在\trIsomorphism{} $\Gamma(X,\mathscr{F})_f\xrightarrow{\ \sim\ }\Gamma(X_f,\mathscr{F})$，使下图交换。}
\[
\begin{tikzcd}[row sep=2.6em, column sep=5.5em]
\Gamma(X,\mathscr{F}) \ar[rr,"\mathrm{res}_{X_f\subset X}"] \ar[dr,"{\otimes_{\Gamma(X,\OO_X)}\bigl(\Gamma(X,\OO_X)_f\bigr)}"'] && \Gamma(X_f,\mathscr{F}) \\
& \Gamma(X,\mathscr{F})_f \ar[ur,"\sim"] &
\end{tikzcd}
\]

\textbf{6.2.G.（重要练习）}\quad 证明 Qcqs 引理~6.2.9。提示：「拟紧拟分离」这条假设恰好翻译成：「被有限多个\trAffine\trOpenSet{}所覆盖，且它们两两的交也被有限多个\trAffine\trOpenSet{}所覆盖」。把\trExact\trFunctor{} $\otimes_{\Gamma(X,\OO_X)}\Gamma(X,\OO_X)_f$ 作用到\trExactSequence
\[
0\longrightarrow\Gamma(X,\mathscr{F})\longrightarrow\bigoplus_i\Gamma(\mathscr{U}_i,\mathscr{F})\longrightarrow\bigoplus\Gamma(\mathscr{U}_{ijk},\mathscr{F}),
\]
之上，其中诸 $\mathscr{U}_i$ 构成 $X$ 的一个有限仿射\trCover，而诸 $\mathscr{U}_{ijk}$ 构成 $\mathscr{U}_i\cap\mathscr{U}_j$ 的一个有限仿射\trCover。

\textbf{6.2.H.（不重要练习）}\quad 给出一个反例，用以说明：没有\trQuasicompactness{}假设，Qcqs 引理~6.2.9 便不成立。（可能的提示：取\trAffineScheme{}的一个无穷\trDisjointUnion。关键想法是：无穷积与\trLocalization{}不交换。）

\textbf{6.2.10. \extrahard\ \trGrothendieck\trTopology。}\quad \trDistinguishedAffineBase{}并不是通常意义下的\trTopology{}——比如说，两个仿射集的并未必还是仿射集。然而，它确是「\trTopology{}的一种推广」的头一个新例子——即\textbf{\trSite}（site）或\textbf{\trGrothendieck\trTopology}（Grothendieck topology）的观念。我们给出定义，以满足好奇者的好奇心，但我们肯定不会使用这个观念。（想要一份干净的表述，可参看 [Stacks, tag~00VH]；我们此处的讨论只打算充当动机。）其想法是：我们应当抽象出的，仅仅是定义\trSheaf{}所必需的那些观念。我们需要\trOpenSet{}的观念，但结果是我们甚至不需要\trUnderlyingSet，也就是说，我们甚至不需要点的观念！让我们细想一下，我们所需要的东西有多么少。为了使我们对\trSheaf{}的讨论得以运转，我们需要知道\trOpenSet{}是什么，（被允许的）\trInclusionMap{}是什么，以及它们应当「表现良好」；特别地，\trOpenSet{}与\trInclusionMap{}的数据应当构成一个\trCategory{}。（例如，一个被允许的\trInclusionMap{}与另一个被允许的\trInclusionMap{}的\trComposition{}应当仍是被允许的\trInclusionMap——在\trDistinguishedAffineBase{}中，\trDistinguishedOpen{}的\trDistinguishedOpen{}是\trDistinguishedOpen。）于是我们只要求\emph{这个\trCategory}的数据。到这一步，我们已经能够定义\trPresheaf{}了（无非是从这个由「\trOpenSet」组成的\trCategory{}出发的一个\trContravariantFunctor{}）。这个想法我们在\cref{ex:presheaf-contravariant-functor}中早已见过。

为了把这个定义扩充成\trSheaf{}的定义，我们还需要更多的信息。我们希望一个\trOpenSet{}的两个\trOpenSubset{}相交成一个\trOpenSet，所以我们\emph{希望这个\trCategory{}对\trFiberProduct{}封闭}（对照\cref{ex:fiber-products-of-open-sets}）。对于\trIdentityAxiom{}与\trGluability{}公理，我们需要知道\emph{一些\trOpenSet{}何时覆盖另一个\trOpenSet}，所以我们把这一点也纳入\trGrothendieck\trTopology{}的数据之中。覆盖所满足的数据具有一些显然的性质。每个\trOpenSet{}都覆盖它自身（也就是说，\emph{\trOpenSet{}的\trCategory{}中的\trIdMorphism{}是一个\trCover}）。\trCover{}可以\trPullback：\emph{若有映射 $Y\to X$，则 $X$ 的任何\trCover{}都\trPullback{}成 $Y$ 的一个\trCover}。最后，\trCover{}的\trCover{}\emph{应当是\trCover}。这样的数据（满足这些公理者）就称为一个\trGrothendieck\trTopology{}或一个\textbf{\trSite}。（文献中还有这个定义的若干有用变体。再说一遍，我们遵循 [Stacks]。）「\trGrothendieck\trTopology{}上的\trSheaf」这一观念，我们可以按惯常的方式来定义，无须任何改动。\textbf{\trTopos}（topos）是给「一个\trGrothendieck\trTopology{}上由\trSet{}的\trSheaf{}组成的\trCategory{}」起的吓人名字。

\trGrothendieck\trTopology{}被用在代数几何及其周边的各式各样情境中。\trEtale{}\trCohomology{}（使用\trEtale{}\trTopology）是\trGalois\trCohomology{}的推广，是一件核心工具；更一般的平坦\trTopology{}也是如此，例如光滑\trTopology。\trDeligne--\trMumford\trStack{}与\trArtin\trStack{}的定义分别用到\trEtale{}\trTopology{}与光滑\trTopology。\trTate{}借助一个合适的\trGrothendieck\trTopology，为诸如 $\mathbb{Q}_p$ 这类\trTotallyDisconnected{}\trBaseField{}上的\trNonArchimedean{}解析几何发展了一套很好的理论。$K$~理论方面的工作（例如与\trVoevodsky{}（Voevodsky）的工作有关系的那些）使用各种奇异的\trTopology。
